% 基数约束 k=2 的边际增益贪心（传感器覆盖实例）。
% 第1轮增益 a10 b8 c8 d1，选a；第2轮（已选{a}）增益 b2 c4 d1，选c；贪心 {a,c} 覆盖14格。
% 枚举全部至多2元子集，最优为 {b,c}，覆盖16格；14 >= (1-(1-1/2)^2)*16 = 12。
\begin{tikzpicture}[x=1mm,y=1mm,line width=.6pt,font=\figurefont\fontsize{8}{10}\selectfont]
  \def\u{1.9}
  % 列标题
  \node[FigureInk] at (12,32) {第1轮：已选$\varnothing$};
  \node[FigureInk] at (44,32) {第2轮：已选$\{a\}$};
  \foreach \v/\y in {a/24,b/18,c/12,d/6}{\node[FigureInk] at (-4,\y) {$\v$};}
  % 第1轮
  \foreach \y/\g/\ch in {24/10/1,18/8/0,12/8/0,6/1/0}{
    \ifnum\ch=1 \fill[FigureYellow] (0,\y-2) rectangle (\g*\u,\y+2);
    \else \fill[FigureBlueFill] (0,\y-2) rectangle (\g*\u,\y+2); \draw[FigureBlue,line width=.6pt] (0,\y-2) rectangle (\g*\u,\y+2); \fi
    \node[FigureInk,anchor=west] at (\g*\u+.8,\y) {\g};
  }
  % 第2轮
  \node[FigureMuted,anchor=west] at (32,24) {已选};
  \foreach \y/\g/\ch in {18/2/0,12/4/1,6/1/0}{
    \ifnum\ch=1 \fill[FigureYellow] (32,\y-2) rectangle (32+\g*\u,\y+2);
    \else \fill[FigureBlueFill] (32,\y-2) rectangle (32+\g*\u,\y+2); \draw[FigureBlue,line width=.6pt] (32,\y-2) rectangle (32+\g*\u,\y+2); \fi
    \node[FigureInk,anchor=west] at (32+\g*\u+.8,\y) {\g};
  }
  \draw[FigureGrid] (0,2)--(0,28); \draw[FigureGrid] (32,2)--(32,28);
  % 右侧：覆盖结果对照
  \begin{scope}[shift={(66,17)},x=4.2mm,y=4.2mm]
    \node[FigureInk,anchor=south west] at (0,2.05) {贪心$\{a,c\}$：14格};
    \foreach \c in {1,...,8}{\foreach \y in {0,1}{
      \pgfmathsetmacro{\cov}{\c>1?1:0}
      \ifdim\cov pt>0pt \filldraw[fill=FigureYellowFill,draw=FigureGrid] (\c-1,\y) rectangle (\c,\y+1);
      \else \filldraw[fill=white,draw=FigureGrid] (\c-1,\y) rectangle (\c,\y+1); \fi}}
    \draw[FigureRed,line width=1pt] (0.15,0.15) rectangle (0.85,1.85);
  \end{scope}
  \begin{scope}[shift={(66,1)},x=4.2mm,y=4.2mm]
    \node[FigureInk,anchor=south west] at (0,2.05) {枚举最优$\{b,c\}$：16格};
    \foreach \c in {1,...,8}{\foreach \y in {0,1}{
      \filldraw[fill=FigureBlueFill,draw=FigureGrid] (\c-1,\y) rectangle (\c,\y+1);}}
  \end{scope}
\end{tikzpicture}
